\documentclass[12pt,a4paper]{article}

\usepackage[utf8]{inputenc}
\usepackage[T1]{fontenc}
\usepackage[spanish, es-nodecimaldot]{babel}
\usepackage{geometry}
\usepackage{graphicx}
\usepackage{float}
\usepackage{booktabs}

\geometry{
    left=1in,
    right=1in,
    top=1in,
    bottom=1in
}

\begin{document}

\begin{table}[H]
\centering
\caption{Direccionamiento IPv4 de la topología.}
\label{tab:direccionamiento}
\small

\begin{tabular}{llll}
\toprule
Dispositivo & Interfaz & Dirección IP & Red \\
\midrule
R1 & S0/3/0 & 140.10.0.1/24 & 140.10.0.0/24 \\
R1 & S0/3/1 & 52.0.0.1/24 & 52.0.0.0/24 \\
R1 & G0/0 & 206.10.5.1/24 & 206.10.5.0/24 \\
\midrule
R2 & S0/3/1 & 140.10.0.2/24 & 140.10.0.0/24 \\
R2 & S0/3/0 & 140.30.0.1/24 & 140.30.0.0/24 \\
R2 & S0/2/1 & 140.0.0.1/24 & 140.0.0.0/24 \\
R2 & G0/0 & 205.20.3.1/24 & 205.20.3.0/24 \\
\midrule
R3 & S0/3/1 & 140.30.0.2/24 & 140.30.0.0/24 \\
R3 & S0/3/0 & 140.20.0.1/24 & 140.20.0.0/24 \\
\midrule
R4 & S0/3/0 & 140.0.0.2/24 & 140.0.0.0/24 \\
R4 & S0/3/1 & 140.20.0.2/24 & 140.20.0.0/24 \\
R4 & G0/0 & 148.10.144.1/24 & 148.10.144.0/24 \\
R4 & G0/1 & 148.10.176.1/24 & 148.10.176.0/24 \\
\midrule
ISP & S0/3/0 & 52.0.0.2/24 & 52.0.0.0/24 \\
ISP & G0/0 & 52.0.1.1/24 & 52.0.1.0/24 \\
\bottomrule
\end{tabular}
\end{table}

\end{document}